\section{Obstructions to Further Improvement}\label{sec:obstructions}

This section records three barriers encountered by natural attempts to improve
the upper bound.  They are not used in the proof of \Cref{thm:main-upper}; their
purpose is to delineate which proof classes cannot by themselves resolve the
asymptotic question.

\subsection{A local-relaxation barrier}

The first obstruction says that bounded-round local relaxations can have an
integrality gap at precisely the scale where one would hope for a short
transversal argument.  We record it for sparse $\ell$-uniform families rather
than for the complete hypergraph.

For an $\ell$-uniform family $\mathcal H\subseteq\binom P\ell$, let
$\tau_{\Z}(\mathcal H)$ denote its integral transversal number and let
$\tau_f(\mathcal H)$ denote the value of the fractional set-cover relaxation
\[
  \min \sum_{u\in P}x_u
  \quad\text{subject to}\quad
  \sum_{u\in C}x_u\ge 1\ (C\in\mathcal H),\qquad 0\le x_u\le 1.
\]

\begin{proposition}[Sherali--Adams transversal barrier]\label{prop:sa-barrier}
Fix $0<\alpha<1$.  For all sufficiently large $N$, put
\[
  \ell=\lfloor N/\log N\rfloor,\qquad q=\lfloor \alpha N\rfloor,\qquad
  \delta=\frac{\binom{N-q}{\ell}}{\binom N\ell}.
\]
For every $N$-point set $P$ there is a family
$\mathcal C\subseteq\binom P\ell$ such that every
$Q\subseteq P$ with $|Q|\le q$ is disjoint from some $C\in\mathcal C$ and
\[
  |\mathcal C|\le
  \left\lceil\frac{\log\binom Nq+1}{\delta}\right\rceil .
\]
Consequently
\[
  \tau_{\Z}(\mathcal C)\ge q+1>\alpha N,
  \qquad
  \tau_f(\mathcal C)\le \frac N\ell=(1+o(1))\log N.
\]
Moreover, for every $\ell$-uniform family
$\mathcal H\subseteq\binom P\ell$ and every integer $0\le r<\ell$, the
level-$r$ Sherali--Adams relaxation of the above transversal LP has value at
most
\[
  \frac{N}{\ell-r}.
\]
\end{proposition}

\begin{proof}
Choose
\[
  s=\left\lceil\frac{\log\binom Nq+1}{\delta}\right\rceil
\]
independent uniformly random $\ell$-subsets of $P$.  For a fixed $q$-set
$Q\subseteq P$, a random $\ell$-set avoids $Q$ with probability $\delta$.
Thus the probability that no chosen set avoids this $Q$ is at most
$\exp(-s\delta)\le e^{-1}/\binom Nq$.  A union bound over all $q$-sets gives a
positive probability that every $q$-set is avoided by at least one chosen
$\ell$-set.  Removing duplicate choices gives the desired family
$\mathcal C$.  If $Q'\subseteq P$ has $|Q'|\le q$, extend it to a $q$-set
$Q$; any $C\in\mathcal C$ disjoint from $Q$ is also disjoint from $Q'$.

The avoidance property forces every transversal to have size at least
$q+1$, proving $\tau_{\Z}(\mathcal C)> \alpha N$.  The uniform fractional
assignment $x_u=1/\ell$ satisfies every covering constraint, so
$\tau_f(\mathcal C)\le N/\ell$.

It remains to prove the Sherali--Adams bound.  Let
$\mathcal H\subseteq\binom P\ell$ be arbitrary and fix $0\le r<\ell$.  Set
\[
  p=\frac1{\ell-r}.
\]
For a covering row $C\in\mathcal H$ and disjoint sets
$I,J\subseteq P$ with $|I\cup J|\le r$, write
\[
  L_{I,J}(x)=\prod_{i\in I}x_i\prod_{j\in J}(1-x_j).
\]
In the row-wise lifted formulation, introduce the literal-moment value
\[
  y_{C,I,J}:=\PP(I\subseteq Y,\ J\cap Y=\emptyset)
  =p^{|I|}(1-p)^{|J|},
\]
where $Y$ is a product Bernoulli-$p$ random subset of $P$.  The value assigned
to any lifted monomial $x_uL_{I,J}$ is the corresponding product-measure
moment: it is $0$ if $u\in J$, $y_{C,I,J}$ if $u\in I$, and
$p\,y_{C,I,J}$ if $u\notin I\cup J$.  Since these are moments of an actual
probability distribution on $\{0,1\}^P$, the lifted box constraints are
satisfied; in particular $0\le y_{C,I,J}\le 1$ for every lifted variable.

Now multiply a covering constraint
\[
  g_C(x):=\sum_{u\in C}x_u-1\ge 0
\]
by any literal $L_{I,J}$ with $|I\cup J|\le r$.  If $I\cap J\ne\emptyset$,
then $L_{I,J}\equiv 0$.  Assume $I\cap J=\emptyset$.  If
$I\cap C\ne\emptyset$, then whenever $L_{I,J}=1$ at least one variable in
$C$ is already fixed to $1$, so $g_C(x)L_{I,J}(x)\ge 0$ pointwise.

It remains to handle $I\cap C=\emptyset$.  Let $c=|J\cap C|$.  Conditional on
the event $L_{I,J}=1$, the $c$ variables in $C\cap J$ are fixed to $0$, while
the remaining $\ell-c$ variables of $C$ remain independent Bernoulli-$p$.
Hence
\[
  \mathbb{E}\!\left[\sum_{u\in C}x_u\,\middle|\, L_{I,J}=1\right]
  =(\ell-c)p.
\]
Since $c\le |J|\le |I\cup J|\le r$, we have
\[
  (\ell-c)p\ge (\ell-r)\frac1{\ell-r}=1.
\]
Therefore $\mathbb{E}[g_C(x)L_{I,J}(x)]\ge 0$ for every lifted covering row, so the
product Bernoulli moment vector is feasible for the level-$r$
Sherali--Adams relaxation.  Its objective value is
\[
  \mathbb{E}\sum_{u\in P}x_u=Np=\frac{N}{\ell-r}.
\]
\end{proof}

Thus no argument whose entire certificate is bounded by the above level-$r$
Sherali--Adams relaxation can certify a transversal lower bound exceeding
$N/(\ell-r)$ on these instances, even though the true integral optimum is
$>\alpha N$.  The artifact script
\texttt{scripts/sa\_barrier\_verification.py} exhaustively checks the lifted
covering inequalities for $N\le 7$ using exact rational arithmetic.

\subsection{A q-shadow dichotomy}

The next obstruction is an extremal statement for Johnson shadows.  It explains
why a sparse separator family cannot be forced by expansion alone unless a
large pre-existing shadow is already present.

Let $P$ be a $K$-point set, write $K=h+L$ with $L\ge 1$, and fix
$1\le q\le h$.  Put $X=\binom Pq$ and $Y=\binom Ph$.  A $q$-set $Q$ captures an
$h$-set $H$ if $Q\subseteq H$.  Let $R\subseteq Y$ be a live packet of density
$r:=|R|/|Y|>0$.  Let $D\subseteq X$ be a previously forbidden $q$-shadow, and
let $\mathcal B\subseteq Y$ be a family of already played $h$-sets.  Define
\[
  \delta_q:=\frac{\binom{K-q}{h-q}}{\binom{K}{h}},
\]
the density of $h$-sets containing a fixed $q$-set.  Write
\[
  \partial_h D:=\{H\in Y:\text{some }Q\in D\text{ satisfies }Q\subseteq H\},
  \qquad
  \sigma_q(D):=\frac{|\partial_h D|}{|Y|}.
\]
Thus $\sigma_q(D)$ is an $h$-level shadow measure, not a $q$-level count.  The
normalized second singular value of the $q$-versus-$h$ inclusion graph
satisfies
\[
  \lambda_q^2=\frac{q(K-h)}{h(K-q)}.
\]

\begin{proposition}[q-shadow covering dichotomy]\label{prop:q-shadow}\label{prop:q-shadow-dichotomy}
Let
\[
  U_q(D,\mathcal B):=
  D\cup\bigcup_{B\in\mathcal B}\binom Bq,
  \qquad
  A:=X\setminus U_q(D,\mathcal B)
\]
be the excluded and legal $q$-families.  Assume no legal $q$-separator
$Q\in A$ captures at least
\[
  \frac12\,\delta_q |R|
\]
members of $R$.  Then
\[
  \frac{|U_q(D,\mathcal B)|}{|X|}
  >
  1-\frac{4\lambda_q^2}{r}.
\]
Consequently
\[
  \sigma_q(D)+|\mathcal B|\delta_q
  >
  1-\frac{4\lambda_q^2}{r}.
\]
In particular, if $r\ge\eta>0$, the right-hand side may be weakened to
$1-4\lambda_q^2/\eta$.
\end{proposition}

\begin{proof}
Put $a=|A|/|X|$.  The case $A=\emptyset$ is immediate, so assume $a>0$.
The expander-mixing inequality for the inclusion graph between $X$ and $Y$
gives
\[
  \left|
  e(A,R)-\frac{|A|}{|X|}\frac{|R|}{|Y|}e(X,Y)
  \right|
  \le
  \lambda_q\,e(X,Y)\sqrt{\frac{|A|}{|X|}\frac{|R|}{|Y|}}.
\]
If
\[
  a\ge \frac{4\lambda_q^2}{r},
\]
then the average number of elements of $R$ captured by a member of $A$ is at
least half the random expectation, namely at least
$\frac12\delta_q |R|$.  This contradicts the hypothesis that no legal separator
captures that many live packet members.  Hence
\[
  a<\frac{4\lambda_q^2}{r},
\]
which is exactly
\[
  \frac{|U_q(D,\mathcal B)|}{|X|}
  =1-a
  >
  1-\frac{4\lambda_q^2}{r}.
\]

It remains only to convert this exact $q$-level covering statement into the
displayed shadow bound.  By the LYM normalized matching property on the Boolean
lattice \cite[Section~2.2]{Anderson87}, every
$\mathcal F\subseteq\binom Pq$ satisfies
\[
  \frac{|\partial_h\mathcal F|}{\binom Kh}
  \ge
  \frac{|\mathcal F|}{\binom Kq}
  \qquad(q\le h).
\]
Applying this to $\mathcal F=D$ gives $|D|/|X|\le\sigma_q(D)$.  Also each
$B\in\mathcal B$ contains exactly $\binom hq$ members of $X$, so
$\binom hq/\binom Kq=\delta_q$.  The union bound gives
\[
  \frac{|U_q(D,\mathcal B)|}{|X|}
  \le
  \frac{|D|}{|X|}+|\mathcal B|\delta_q
  \le
  \sigma_q(D)+|\mathcal B|\delta_q,
\]
and the conclusion follows.
\end{proof}

In the central regime relevant to top-facet packet arguments, take
$L\sim h/\log h$ and $q\sim 2(\log h)^2$.  Then
\[
  \delta_q=(e+o(1))h^{-2},
  \qquad
  \lambda_q^2=(2+o(1))\frac{\log h}{h}.
\]
Thus, when $r$ is bounded below and $\sigma_q(D)=o(1)$, separator starvation
forces $|\mathcal B|>(e^{-1}-o(1))h^2$.  Pure expansion cannot supply an
unlimited sequence of fresh separators.

\subsection{A separator-only limitation}

The last obstruction is a finite model for strategies which only play odd
carrier separators and ignore dyadic shielding.

\begin{proposition}[Separator-only limitation]\label{prop:separator-only}
Fix a finite set $P$ of odd primes and a family $\mathcal S\subseteq 2^P$ of
odd squarefree carrier supports.  Suppose Shortener's planned moves are only
odd squarefree separators with supports in $\mathcal S$.  Then there are
dyadic upper-half targets of the form
\[
  2^a\prod_{p\in T}p
\]
which remain untouched by all separators unless $\mathcal S$ covers the
corresponding support $T$ itself.  In particular, any separator-only proof that
does not control all dyadic lifts of the carrier family leaves a positive
residual packet.
\end{proposition}

\begin{proof}
An odd separator can only be comparable with an upper-half target
$2^a m$ through its odd part $m$.  The dyadic exponent is invisible to the odd
separator.  Therefore, if a support $T$ is not covered by the chosen odd
separator family, every dyadic lift $2^a\prod_{p\in T}p$ with value in
$(n/2,n]$ remains incomparable with all separator moves.

For any fixed odd part $m\le n/2$, there is at least one dyadic lift
$2^a m\in(n/2,n]$ unless $m=1$; choose $a$ maximal with $2^a m\le n$.  If
$2^a m\le n/2$, then $2^{a+1}m\le n$, contradicting maximality.  Thus each
uncovered odd support produces an upper-half target not hit by the
separator-only family.  Summing over an uncovered packet gives the claimed
residual.
\end{proof}

The proposition is deliberately modest: it does not rule out separator methods
which also track dyadic lifts or exact upper targets.  It rules out only the
separator-only abstraction in which odd carrier supports are treated as if they
were the whole divisibility geometry.
